\chapter{Orders as integral arithmetic structures}
\label{ch:orders}

\section{The role of the two volumes}

The first attached volume is \emph{Lattices over Orders I}, by Klaus W.
Roggenkamp and Verena Huber-Dyson, published in 1970 as volume 115 of
\emph{Lecture Notes in Mathematics}. Its five chapters develop modules,
homological algebra, Morita equivalence and separability, maximal orders,
and the Higman ideal with extensions of lattices. The second volume is
by Roggenkamp. The first volume is particularly useful here because it
explains precisely which integral information disappears upon extending
scalars to a field. This is the information that must be retained when
one passes between elliptic curves, ideals, finite modules, and geometric
realizations of arithmetic lattices. Our notation uses ordinary
left-to-right function application, whereas the book often places
homomorphisms opposite the scalars. This change of convention must be
remembered when comparing endomorphism rings and opposite rings
\cite{RoggenkampHuberDyson1970}.

The corrections printed at the end of Volume II are part of the
source, not optional editorial detail. They add a restriction
excluding totally definite quaternion matrix components to a
cancellation assertion in the introduction, and replace an invalid
torsion-decomposition step in the proof at chapter-page IV.13 by
an argument using completion and evaluation maps. We use those
corrections when interpreting the structural results; the elementary
proofs below do not depend on the faulty decomposition
\cite[Corrections Vol.~I]{v2Roggenkamp1970II}.

There is a genuine common language, but not a common theorem that
automatically solves all of the problems in this monograph. An order is
an integral algebra. A lattice over it is a module satisfying a finiteness
and torsion condition. A Euclidean lattice carries additional metric
data. A lattice in a lattice gauge theory is a discretization of space
or spacetime. These objects can interact through explicit constructions,
but the shared word ``lattice'' supplies no such construction by itself.
The present chapter builds the arithmetic objects and specifies the
extra structure required at each subsequent step.

The scope of the proofs is deliberate. Elementary statements used in
the later bridges are proved below. Structural theorems whose proofs
require a substantial additional theory are identified as external
results. A proof over an arbitrary field, a theorem over a Dedekind
domain, and a theorem about a positive Hilbert-space operator have
different hypotheses; replacing one of those settings by another is
itself a mathematical obligation.

\section{Full lattices and orders}

Throughout this chapter, let $R$ be a Dedekind domain with fraction
field $K$, and let $A$ be a finite-dimensional separable $K$-algebra.
An $R$-algebra structure on $A$ is understood through the central
embedding $R\subseteq K\subseteq A$. The number-field example is
$R=\mathcal O_K$, but many elementary arguments need only that $R$ be
a noetherian domain. Unless explicitly stated otherwise, modules are
unital and finitely generated.

\begin{definition}[Arithmetic lattice and order]
An $R$-lattice is a finitely generated torsion-free $R$-module. If $V$
is a finite-dimensional $K$-vector space, a full lattice in $V$ is an
$R$-submodule $L\subseteq V$ satisfying $K L=V$ and the preceding
finiteness condition. An $R$-order in $A$ is a subring
$\Lambda\subseteq A$ containing $1_A$, which is a full $R$-lattice
in $A$. A left $\Lambda$-lattice is a left $\Lambda$-module whose
underlying $R$-module is an $R$-lattice.
\end{definition}

The word ``full'' refers to the specified ambient vector space. It
does not mean that the lattice is free over $\Lambda$. Even its
underlying $R$-module need not be free when $R$ is not a principal
ideal domain. The basic Dedekind-domain theorem says that a finitely
generated torsion-free $R$-module is projective. Projectivity over $R$
does not imply projectivity over an arbitrary order $\Lambda$.
This distinction is central to the first volume's treatment of
lattices and extensions \cite[Chs.~I and IV]{RoggenkampHuberDyson1970}.

\begin{proposition}[Common denominators]
\label{prop:common-denominators}
Let $L$ and $M$ be full $R$-lattices in the same finite-dimensional
$K$-space $V$. There exist nonzero $r,s\in R$ such that
$rL\subseteq M$ and $sM\subseteq L$. Consequently $L\cap M$ is
full, and both $(L+M)/L$ and $(L+M)/M$ are finitely generated torsion
$R$-modules.
\end{proposition}

\begin{proof}
Choose finite generating sets $l_1,\ldots,l_a$ of $L$ and
$m_1,\ldots,m_b$ of $M$. Since $KM=V$, each $l_i$ is a finite
$K$-linear combination of the $m_j$. A common nonzero denominator
for these finitely many coefficients gives $r\in R\setminus\{0\}$
with $r l_i\in M$ for every $i$, hence $rL\subseteq M$.
Interchanging the two lattices gives $s$. Thus $rL$ lies in
$L\cap M$ and spans $V$, while $L\cap M$ is finitely generated
because it is a submodule of the noetherian module $L$. The quotient
$(L+M)/L$ is finitely generated and killed by $s$; the other quotient
is killed by $r$.
\end{proof}

Over the ring of integers of a number field, a finitely generated
torsion module is finite as a set. Over an arbitrary Dedekind domain,
that last conclusion requires finite residue fields and must not be
silently inserted. Even in the number-field case, finiteness gives no
useful complexity bound until the denominators and their bit lengths
are controlled.

\begin{proposition}[A stable integral model]
\label{prop:stable-model}
If $V$ is a finite-dimensional left $A$-module and $\Lambda$ is an
$R$-order in $A$, then $V$ contains a full left $\Lambda$-lattice.
\end{proposition}

\begin{proof}
Choose a $K$-basis $v_1,\ldots,v_n$ of $V$ and put
$L=\sum_{i=1}^n\Lambda v_i$. Multiplication by $\Lambda$ preserves
$L$. If $a_1,\ldots,a_t$ generate $\Lambda$ over $R$, then the
finitely many vectors $a_jv_i$ generate $L$ over $R$. As a submodule
of a $K$-vector space, $L$ is torsion-free. Since $1\in\Lambda$,
the vectors $v_i$ lie in $L$, so $KL=V$.
\end{proof}

This simple construction explains why rational representations admit
integral models, while also explaining their nonuniqueness. Choosing a
basis is not the same as classifying all stable lattices. The scalar
extension functor $K\otimes_R-$ can send nonisomorphic lattices to
isomorphic $A$-modules.

\begin{example}[An order that retains a congruence]
For a nonzero ideal $\mathfrak a\subseteq R$, consider
\[
 \Lambda_{\mathfrak a}=
 \left\{\begin{pmatrix}a&b\\c&d\end{pmatrix}:
 a,b,d\in R,\ c\in\mathfrak a\right\}\subseteq M_2(K).
\]
The lower-left entry of a product is $ca'+dc'$, which again belongs
to $\mathfrak a$. Thus this is an order. Tensoring with $K$ removes
the condition $c\in\mathfrak a$, giving all of $M_2(K)$. The
congruence data remains relevant to integral module theory even
though the ambient rational algebra is unchanged.
\end{example}

\section{Endomorphism orders and faithful models}

\begin{proposition}[Endomorphisms of a full lattice]
\label{prop:end-order}
Let $L$ be a full $R$-lattice in $V$. Each $R$-linear endomorphism
of $L$ extends uniquely to a $K$-linear endomorphism of $V$, and
\[
 \operatorname{End}_R(L)=
 \{T\in\operatorname{End}_K(V):T(L)\subseteq L\}
\]
is an $R$-order in $\operatorname{End}_K(V)$.
\end{proposition}

\begin{proof}
Torsion-freeness identifies $L$ with its image in $K\otimes_RL=V$.
Tensoring an endomorphism with $K$ supplies the extension; uniqueness
follows because $L$ spans $V$. The displayed set is closed under
addition and composition and contains the identity. Since $L$ is a
finite projective $R$-module, it is a summand of $R^N$ for some $N$.
Its endomorphism module is therefore a summand of, in particular a
finite submodule of, $\operatorname{Hom}_R(R^N,R^N)$, so it is
finitely generated and torsion-free. Finally, if
$T\in\operatorname{End}_K(V)$, common denominators for the images
of finitely many generators of $L$ give nonzero $r\in R$ with
$rT(L)\subseteq L$. Hence every $T$ belongs to
$K\operatorname{End}_R(L)$.
\end{proof}

For a full lattice $I\subseteq A$, its two multiplier orders are
\[
 O_L(I)=\{a\in A:aI\subseteq I\},\qquad
 O_R(I)=\{a\in A:Ia\subseteq I\}.
\]
The distinction matters in a quaternion algebra: a left ideal of
$O$ usually has a different right order. A computation that changes
the side without changing the convention can reverse composition or
identify the wrong endomorphism ring.

\begin{proposition}[Multiplier orders]
Both $O_L(I)$ and $O_R(I)$ are $R$-orders in $A$.
\end{proposition}

\begin{proof}
Consider $O_L(I)$. Left multiplication embeds $A$ into
$\operatorname{End}_K(A)$, since an element acting as zero sends
$1$ to zero. The inverse image of $\operatorname{End}_R(I)$ is
therefore a torsion-free $R$-module, and it is finite because it is
isomorphic to a submodule of the finite module
$\operatorname{End}_R(I)$. It is a ring containing $1$. To prove
fullness, choose any full lattice $L\subseteq A$. For finite
generators of $L$ and $I$, clear the finitely many denominators
needed to put each product into $I$. This gives nonzero $r\in R$
such that $(rL)I\subseteq I$, so $rL\subseteq O_L(I)$. The same
argument using right multiplication proves the other assertion.
\end{proof}

An arithmetic representation can now be studied through three
related objects: its module $L$, the order acting on the left, and
the commuting endomorphisms acting on the right. The first volume
systematizes this viewpoint through generators and Morita theory
\cite[Ch.~III]{RoggenkampHuberDyson1970}. One must nevertheless keep
the integral module, not just its rational endomorphism algebra.
For example, knowing that the latter is a quadratic field does not
determine the conductor of the integral order.

\section{Localization, support, and what local data detects}

For a maximal ideal $\mathfrak p$ of $R$, write
$R_{\mathfrak p}$ for the localization and
$M_{\mathfrak p}=R_{\mathfrak p}\otimes_R M$. Write
$\widehat R_{\mathfrak p}$ for the completion. These are distinct
constructions. A statement proved for complete discrete valuation
rings does not automatically hold for an uncompleted local ring.
The first volume treats localization and completion separately
\cite[Ch.~I, \S\S6--9]{RoggenkampHuberDyson1970}.

\begin{lemma}[A local membership test]
\label{lem:local-membership}
Let $M$ be an $R$-module and $x\in M$. If the image of $x$ in
$M_{\mathfrak p}$ is zero for every maximal ideal $\mathfrak p$,
then $x=0$.
\end{lemma}

\begin{proof}
Suppose $x\ne0$. Its annihilator
$J=\{r\in R:rx=0\}$ is a proper ideal, hence is contained in
a maximal ideal $\mathfrak p$. If $x/1=0$ in
$M_{\mathfrak p}$, the definition of localization supplies
$s\notin\mathfrak p$ with $sx=0$. This puts $s\in J\subseteq
\mathfrak p$, a contradiction.
\end{proof}

\begin{proposition}[Equality from localized inclusions]
If $L,M\subseteq V$ are lattices, then $L\subseteq M$ if and
only if $L_{\mathfrak p}\subseteq M_{\mathfrak p}$ for every
maximal ideal $\mathfrak p$. In particular, equality of the
embedded local lattices implies $L=M$.
\end{proposition}

\begin{proof}
One implication follows from functoriality. For the converse,
apply Lemma~\ref{lem:local-membership} to the image of each
$l\in L$ in $(L+M)/M$. Localization is exact, so all its
localized images vanish.
\end{proof}

Equality of embedded lattices is stronger than abstract local
isomorphism. To see why, let $R$ be a Dedekind domain with a
nonprincipal invertible ideal $I$. For each $\mathfrak p$, the
local ring $R_{\mathfrak p}$ is a discrete valuation ring and
$I_{\mathfrak p}$ is a free module of rank one. Thus
$I_{\mathfrak p}\simeq R_{\mathfrak p}$ for every prime. But
$I\not\simeq R$ as an $R$-module: an isomorphism would send
$1$ to a generator of $I$, making it principal. The missing data
is the ideal class, an obstruction to making the local
trivializations compatible globally.

\begin{definition}[Genus]
Two $\Lambda$-lattices are in the same genus if their completions
are isomorphic as $\widehat\Lambda_{\mathfrak p}$-modules for
every nonzero prime $\mathfrak p$. When a local-isomorphism version
is used, its equivalence to this definition must be established in
the applicable setting.
\end{definition}

Roggenkamp's 1969 work on irreducible lattices studies precisely the
interaction between maximal overorders, commuting endomorphism
orders, and classes within a genus. Its opening proposition also
supplies a useful cancellation-of-denominators observation: if two
lattices already carry the action of an overorder, a homomorphism
linear over the smaller order is linear over the larger one
\cite[Prop.~1]{Roggenkamp1969Irreducible}. We prove that observation
in our conventions because it prevents several side-changing errors.

\begin{lemma}[Linearity over an overorder]
Let $\Lambda\subseteq\Gamma$ be $R$-orders in $A$, and let
$M,N$ be $\Gamma$-lattices. Every $\Lambda$-linear map
$f:M\to N$ is $\Gamma$-linear.
\end{lemma}

\begin{proof}
By Proposition~\ref{prop:common-denominators}, there is nonzero
$r\in R$ with $r\Gamma\subseteq\Lambda$. For
$\gamma\in\Gamma$ and $m\in M$, $R$-linearity and
$\Lambda$-linearity give
\[
 r f(\gamma m)=f((r\gamma)m)=(r\gamma)f(m)
             =r\gamma f(m).
\]
Torsion-freeness of $N$ permits cancellation of $r$, giving
$f(\gamma m)=\gamma f(m)$.
\end{proof}

The hypothesis that both modules carry the overorder action is
essential. It does not assert that every lattice over $\Lambda$
extends to a lattice over every maximal order containing $\Lambda$.
Extending the action is itself a structural problem.

\section{An explicit Morita equivalence}

An equivalence of module categories is one of the strongest algebraic
bridges available in order theory. It preserves module-theoretic
information, but it does not preserve a metric that was never part
of the category. The following elementary case makes both facts
visible without invoking the general Morita theorem.

\begin{theorem}[Matrix Morita equivalence]
\label{thm:matrix-morita}
Let $S$ be a unital ring, $B=M_n(S)$ with $n\ge1$, and
$e=E_{11}$. The functors
\[
 F(M)=eM,\qquad G(N)=S^n\otimes_S N
\]
give inverse equivalences between left $B$-modules and left
$S$-modules. Here $eBe\simeq S$, and $S^n$ is the usual
$(B,S)$-bimodule of columns.
\end{theorem}

\begin{proof}
For a left $B$-module $M$, define
\[
 \alpha_M:S^n\otimes_S eM\longrightarrow M,
 \qquad e_i\otimes m\longmapsto E_{i1}m.
\]
The balancing relation is respected because multiplication by an
entry of $S$ in a column agrees with the action of the corresponding
corner matrix on $eM$. An inverse is
\[
 \beta_M(m)=\sum_{i=1}^n e_i\otimes E_{1i}m.
\]
Each $E_{1i}m$ lies in $eM$. The identities
$\sum_iE_{i1}E_{1i}=\sum_iE_{ii}=1$ give
$\alpha_M\beta_M=1_M$. On a generator $e_j\otimes m$,
$\beta_M\alpha_M$ equals
$\sum_i e_i\otimes E_{1i}E_{j1}m=e_j\otimes m$,
since $E_{11}m=m$. The maps commute with module homomorphisms and
$\alpha_M$ is $B$-linear, as is checked on matrix units. Finally,
$e(S^n\otimes_SN)$ consists of the first-coordinate copy of $N$,
giving a natural isomorphism $FG(N)\simeq N$. These natural
isomorphisms prove the equivalence.
\end{proof}

If $S$ is an $R$-order, these functors restrict to the appropriate
finite $R$-lattices: taking a direct summand preserves finite
generation and torsion-freeness, and $G$ is a finite direct sum on
underlying $R$-modules. More general Morita equivalences replace
$S^n$ by a finitely generated projective generator, or
\emph{progenerator}; their treatment is a central component of
Roggenkamp and Huber-Dyson's first volume
\cite[Ch.~III, \S\S1--2]{RoggenkampHuberDyson1970}.

To connect this with complexes, apply $F$ degree by degree. Since
$F$ is exact, kernels and images are preserved up to canonical
isomorphism, and therefore
\[
 H^i(F(C^\bullet))\simeq F(H^i(C^\bullet)).
\]
This is a categorical statement about homology. It provides no
identity between spectra of analytic Laplacians placed on the two
complexes after choosing unrelated inner products.

\begin{example}[The missing metric]
The abelian groups $\mathbb Z^2$ with quadratic forms
$Q_1(x,y)=x^2+y^2$ and
$Q_T(x,y)=T^2x^2+T^{-2}y^2$, $T>0$, are identical as
$\mathbb Z$-modules. Their covolumes in the corresponding
Euclidean realizations are both one, but their shortest nonzero
lengths are $1$ and $\min(T,T^{-1})$. No theorem of module
equivalence can force those numbers to agree. Any proposed
cryptographic or spectral transfer must additionally control the
chosen form and the distortion of the transporting maps.
\end{example}

\section{Trace pairings, dual lattices, and discriminants}

Let $t:A\to K$ be a $K$-linear functional such that
$\langle x,y\rangle=t(xy)$ is nondegenerate and
$t(xy)=t(yx)$. Important examples are the field trace in a
separable field extension and the reduced trace in a central simple
algebra. We explicitly assume $t(\Lambda)\subseteq R$ whenever
integrality of a pairing on $\Lambda$ is used. The distinction
between an ordinary matrix trace and a reduced trace is relevant
in positive characteristic.

\begin{definition}[Trace dual]
For a full lattice $L\subseteq A$, define
\[
 L^{\#}=\{x\in A:t(xL)\subseteq R\}.
\]
When $R=\mathbb Z$, $A=K$ is a number field, and
$L=\mathcal O_K$, the trace dual is the inverse different,
also called the codifferent.
\end{definition}

\begin{proposition}[Trace duality]
\label{prop:trace-dual}
The lattice $L^{\#}$ is naturally isomorphic to
$\operatorname{Hom}_R(L,R)$ under $x\mapsto(y\mapsto t(xy))$.
It is full and $(L^{\#})^{\#}=L$.
\end{proposition}

\begin{proof}
Nondegeneracy identifies $A$ with $\operatorname{Hom}_K(A,K)$.
Every $R$-linear functional $L\to R$ extends uniquely to a
$K$-linear functional on $KL=A$, so the displayed map is a
bijection onto $\operatorname{Hom}_R(L,R)$. The latter is a
finite projective module of the same rank as $L$, proving fullness.
Since finite projective modules are canonically isomorphic to their
double duals, applying the identification twice gives the final
equality inside $A$. The elementary double-dual statement follows
first for $R^n$ by coordinates and then for its direct summands.
\end{proof}

There is a bimodule structure compatible with this identification:
if $\phi\in\operatorname{Hom}_R(\Lambda,R)$, put
$(a\phi b)(x)=\phi(bxa)$. Cyclicity of the trace converts this
into left and right multiplication on $\Lambda^{\#}$.
The trace dual need not be a ring. Treating it as an overorder
without checking multiplicative closure would confuse a dual
module with an algebra.

When $R=\mathbb Z$ and $L$ has basis $b_1,\ldots,b_d$, its trace
Gram matrix is $G=(t(b_ib_j))_{ij}$. In these coordinates,
$L^{\#}$ has coordinate module $G^{-1}\mathbb Z^d$; here
symmetry of $G$ follows from the assumption on $t$. If the trace
pairing is integral on $L$, then
\[
 [L^{\#}:L]=|\det G|.
\]
This equality follows by multiplying coordinates by $G$: the
quotient is isomorphic to $\mathbb Z^d/G\mathbb Z^d$, whose
order is the absolute determinant. The determinant formula is
understood for a nondegenerate pairing; a singular Gram matrix
does not define a finite dual quotient.

\begin{proposition}[Change of lattice and discriminant]
Let $L\subseteq M$ be full $\mathbb Z$-lattices in $A$ of finite
index $h$. Then
\[
 \operatorname{disc}_t(L)=h^2\operatorname{disc}_t(M).
\]
\end{proposition}

\begin{proof}
If $C\in M_d(\mathbb Z)$ expresses a basis of $L$ in a basis
of $M$, their Gram matrices satisfy $G_L=C^{\mathsf T}G_MC$.
Taking determinants gives
$\det G_L=(\det C)^2\det G_M$. The index is
$|\det C|$, yielding the formula.
\end{proof}

This is one of the cleanest quantitative bridges in the subject.
The underlying change of integral structure is measured by an
index; the bilinear invariant changes by its square. Yet the
discriminant alone does not determine the shortest vector or the
complexity of a search problem. Different positive forms can share
a determinant and have very different minima.

\section{From arithmetic lattices to Euclidean geometry}

For a number field $K$ with $r_1$ real embeddings and $r_2$ pairs
of complex embeddings, the Minkowski embedding sends $x$ to
\[
 \sigma(x)=\bigl(\sigma_1(x),\ldots,\sigma_{r_1}(x),
 \sqrt2\operatorname{Re}\tau_1(x),
 \sqrt2\operatorname{Im}\tau_1(x),\ldots\bigr)
 \in\mathbb R^{r_1+2r_2}.
\]
The factor $\sqrt2$ is a convention that makes the squared norm
$\sum_i|\sigma_i(x)|^2+2\sum_j|\tau_j(x)|^2$. The real-vector
space isomorphism
$K\otimes_{\mathbb Q}\mathbb R\simeq
\mathbb R^{r_1}\times\mathbb C^{r_2}$ makes the image of any full
$\mathbb Z$-lattice discrete and cocompact. A different coordinate
basis describes the same lattice, whereas a different positive form
changes the geometry.

The field trace $\operatorname{Tr}_{K/\mathbb Q}(xy)$ is
nondegenerate, but it need not be positive definite. In a totally
imaginary quadratic field, $\operatorname{Tr}(x\bar x)$ supplies
a positive form, with a specific use of complex conjugation. In a
general number field, the embedding-based Euclidean form is the
safe uniform construction. Writing a trace expression without
specifying the involution can inadvertently replace a positive
form by an indefinite one.

\begin{proposition}[When a transport controls minima]
Let $T:V\to W$ be a linear isomorphism of finite-dimensional
real inner-product spaces with $\dim V>0$, and suppose
$a\|v\|\le\|Tv\|\le b\|v\|$ for all $v$, with
$0<a\le b$. For a full Euclidean lattice $L\subseteq V$,
\[
 a\lambda_1(L)\le\lambda_1(TL)\le b\lambda_1(L),
\]
where $\lambda_1$ is the length of a shortest nonzero vector.
\end{proposition}

\begin{proof}
The lower bound follows by applying the lower inequality to every
nonzero lattice vector and taking the minimum. The upper bound
follows by applying the upper inequality to a vector realizing
$\lambda_1(L)$. Such a vector exists because a discrete lattice
has only finitely many points in any bounded ball.
\end{proof}

This proof also indicates what a useful algorithmic reduction would
need: a computable map, polynomially bounded bit complexity, and a
controlled ratio $b/a$. An abstract isomorphism of modules provides
none of these quantitative guarantees. They must be proved for the
specific encoding and norm used by a cryptographic problem.

\section{Quaternion orders as four-dimensional arithmetic}

Over a field $K$ of characteristic different from two, for $a,b\in K^\times$ a quaternion
algebra $(a,b)_K$ has basis $1,i,j,ij$ with
$i^2=a$, $j^2=b$, and $ji=-ij$. The standard involution sends
\[
 x=x_0+x_1i+x_2j+x_3ij\quad\hbox{to}\quad
 \bar x=x_0-x_1i-x_2j-x_3ij.
\]
Its reduced trace and norm are
\[
 \operatorname{trd}(x)=x+\bar x=2x_0,\qquad
 \operatorname{nrd}(x)=x\bar x
  =x_0^2-a x_1^2-b x_2^2+ab x_3^2.
\]
The presentation is a convenient coordinate system, not a guarantee
that the displayed $\mathbb Z$-span is maximal when $K=\mathbb Q$.

\begin{proposition}[Multiplicativity of the quaternion norm]
For $x,y$ in a quaternion algebra,
$\operatorname{nrd}(xy)=\operatorname{nrd}(x)
\operatorname{nrd}(y)$.
\end{proposition}

\begin{proof}
The defining relations give
$\overline{xy}=\bar y\bar x$. As $y\bar y$ is a scalar,
\[
 (xy)\overline{xy}=xy\bar y\bar x
 =\operatorname{nrd}(y)x\bar x
 =\operatorname{nrd}(y)\operatorname{nrd}(x).
\]
\end{proof}

For the real Hamilton algebra $\mathbb H=(-1,-1)_{\mathbb R}$,
the reduced norm is the ordinary sum of four squares. A definite
rational quaternion algebra becomes $\mathbb H$ after tensoring
with $\mathbb R$, so a full order is a positive four-dimensional
Euclidean lattice after choosing coordinates. The three-dimensional
trace-zero space is useful for visualizing ternary norm forms. A
three-dimensional projection of the full four-dimensional lattice
can hide distinct points and shorten projected vectors; it must
not be used as a shortest-vector certificate.

Left and right orders of a quaternion ideal describe how its two
ends act. Multiplication of ideals is compatible only when those
ends match. If $I$ is an $(O_1,O_2)$-bimodule and $J$ is an
$(O_2,O_3)$-bimodule, then $I\otimes_{O_2}J$ is an
$(O_1,O_3)$-bimodule. The multiplication map from this tensor
product into the product lattice $IJ$ requires its own hypotheses
to be an isomorphism; it is an isomorphism in the invertible-ideal
setting used for Morita transport.

\section{Maximality, separability, and their limits}

\begin{definition}
An order is maximal if it is not properly contained in any other
$R$-order in the same $K$-algebra. An $R$-algebra $\Lambda$ is
separable if it is projective as a module over
$\Lambda^e=\Lambda\otimes_R\Lambda^{\mathrm{op}}$, with the
action $(a\otimes b^{\mathrm{op}})x=axb$.
\end{definition}

The separability of $A$ over $K$ does not imply the separability of
an integral order over $R$. Ramification is an obstruction.
Roggenkamp's short paper \emph{A Remark on Separable Orders}
characterizes separable orders over rings of integers using
maximality together with unramified conditions on the centers and
the simple algebra components. This result is a useful warning:
``maximal'' and ``separable over the base ring'' cannot be used
interchangeably \cite{Roggenkamp1969Separable}.

For the main arithmetic setting, the first volume develops the
existence and local structure of maximal orders in separable
algebras, and their relationship with progenerators
\cite[Ch.~IV]{RoggenkampHuberDyson1970}. Applying those results to
an algorithm requires a representation of the order, arithmetic in
its coefficient ring, and a procedure for recognizing membership.
The abstract existence theorem does not itself construct a basis
with useful complexity guarantees.

There is also a conceptual boundary worth stating precisely. An
order can support a category of modules, a Grothendieck group of
projectives, and finite complexes. Those are algebraic invariants.
An elliptic operator requires a smooth manifold, vector bundles,
and a principal symbol satisfying an exactness or invertibility
condition away from the zero cotangent vector. A quantum field
theory requires further analytic and probabilistic data. The
intervening structures are the subject of subsequent chapters;
they cannot be inferred merely from the existence of a maximal
order.

\section{A finite certificate for maximal overorders}

There is a useful elementary special case of the existence theory.
Assume now $R=\mathbb Z$ and that $A$ carries the nondegenerate
trace pairing used above, with the additional property that
$t(\Gamma)\subseteq\mathbb Z$ for every order $\Gamma$ in $A$.
This holds for the field trace on number fields and for reduced
trace on rational quaternion algebras. The trace-integrality fact
uses that an element of an order is integral and that the relevant
trace of an integral element is an integer; it is part of the
separable-algebra machinery in the first volume
\cite[III, \S6; IV, Theorem~1.4]{RoggenkampHuberDyson1970}.

\begin{proposition}[Bounding all overorders]
If $\Lambda$ is an order and $\Gamma$ is an order containing it,
then $\Lambda\subseteq\Gamma\subseteq\Lambda^{\#}$. In
particular, there are only finitely many overorders of $\Lambda$,
and at least one of them is maximal.
\end{proposition}

\begin{proof}
For $\gamma\in\Gamma$ and $\lambda\in\Lambda$, their product
belongs to $\Gamma$, so $t(\gamma\lambda)\in\mathbb Z$.
This is exactly the condition $\gamma\in\Lambda^{\#}$.
The quotient $\Lambda^{\#}/\Lambda$ is a finite abelian group.
An intermediate lattice is determined by its subgroup in this
quotient, and a finite group has finitely many subgroups.
Overorders are a subset of those intermediate lattices; hence
they form a finite nonempty partially ordered set. A maximal
element $\Gamma$ of this set cannot have a proper overorder,
because any such overorder would still contain $\Lambda$.
\end{proof}

This argument gives a finite verification strategy for small
examples. Compute a trace-dual basis, enumerate the relevant
finite quotient's subgroups, lift each subgroup to a lattice,
and test that the lifted lattice contains the identity and is
closed under multiplication. A multiplication test requires only
the products of basis elements: bilinearity then covers every
pair of elements. Each product can be expressed exactly in the
candidate basis, and membership is equivalent to integral
coordinates.

The proposition is an existence and certification statement,
not a polynomial-time maximal-order algorithm. The number of
subgroups can be large, and the quotient size depends on the
discriminant. A complexity theorem would have to control those
quantities in terms of the bit length of the input. The general
existence theorem in the book is IV.4.6; local detection of
maximality is IV.4.8, and the hereditary property is IV.4.19
\cite{RoggenkampHuberDyson1970}. These are separate theorems,
not interchangeable definitions.

\section{Averaging and integral extension obstructions}

The first volume's final chapter studies projective homomorphisms
and ideals that annihilate extensions. A finite-group example
shows the mechanism directly and also explains why a denominator
that is harmless over a field can matter over an integral base
\cite[Ch.~V]{RoggenkampHuberDyson1970}.

\begin{proposition}[An integral averaging lift]
Let $G$ be a finite group, let $R$ be a commutative ring, and
let
\[
 0\longrightarrow N\longrightarrow M
   \mathop{\longrightarrow}^{q} L\longrightarrow0
\]
be an exact sequence of left $R[G]$-modules that splits as a
sequence of $R$-modules. Then there is an $R[G]$-linear map
$T:L\to M$ with $qT=|G|\,1_L$.
\end{proposition}

\begin{proof}
Choose an $R$-linear section $s:L\to M$ of $q$, and define
\[
 T(l)=\sum_{g\in G}g\,s(g^{-1}l).
\]
The map is $R$-linear. For $h\in G$, replacing $g$ by $hu$
in the sum gives
\[
 T(hl)=\sum_{u\in G}hu\,s(u^{-1}l)=hT(l),
\]
so $T$ is $R[G]$-linear. Since $q$ is $G$-equivariant and
$qs=1_L$, each summand satisfies
$q(g\,s(g^{-1}l))=g(g^{-1}l)=l$. Summing gives
$qT(l)=|G|l$.
\end{proof}

If $|G|$ is a unit in $R$, dividing $T$ by $|G|$ produces
an equivariant section, so the original exact sequence splits
as $R[G]$-modules. This is the averaging mechanism behind
Maschke's theorem in the field case. If $|G|$ is not invertible,
the proof still supplies a lift of a multiple of the identity,
but not necessarily a section. This residual multiple is
precisely the kind of information measured by annihilators of
extension classes.

For completeness, one can state the extension consequence
without concealing the role of the lift. Pull the sequence back
along the map $|G|1_L:L\to L$. The pullback module is
\[
 P=\{(m,l)\in M\oplus L:q(m)=|G|l\}.
\]
The map $l\mapsto(T(l),l)$ is an $R[G]$-linear section of
$P\to L$. Thus this pullback extension is split. Under the
standard identification of extension classes with
$\operatorname{Ext}^1_{R[G]}(L,N)$, pullback by scalar
multiplication is multiplication by that scalar, so $|G|$
annihilates the extension class. The identification itself is
developed in the book's Chapter II, \S5.

The condition that the sequence split over $R$ holds in the
arithmetic-lattice setting when $L$ is projective as an
$R$-module. It need not hold for arbitrary $R$-modules. The
argument therefore records its exact input rather than
silently extending an assertion about lattices to all modules.

\begin{example}[A denominator that cannot be removed]
Let $k$ have characteristic $p$, and let $G=C_p$ have generator
$g$. The two-dimensional representation with
$g\mapsto\left(\begin{smallmatrix}1&1\\0&1\end{smallmatrix}\right)$
is well-defined because $(1+N)^p=1$ when $N^2=0$.
Its first-coordinate line is a trivial subrepresentation,
and the quotient is also trivial. There is no invariant
complement: a vector $(a,1)$ spanning a candidate complement
is sent to $(a+1,1)$, which is not in its span. Thus the
extension is nonsplit. Here $|G|=p=0$ in $k$, so the
averaging proposition does not supply a nonzero multiple of
the identity that can be inverted.
\end{example}

This example concerns an exact sequence and a representation.
It does not assert a physical mass gap, nor a cryptographic
hardness gap. Its contribution to the broader program is a
precise obstruction: a rational or generic-fiber splitting
can fail integrally at primes dividing a denominator. Such
failure must be retained in any proposed bridge that uses
integral complexes or reductions modulo primes.

\chapter{Elliptic curves through their endomorphism lattices}
\label{ch:elliptic-orders}

\section{Two arithmetic questions attached to the same curve}

An elliptic curve has a group of points and a ring of endomorphisms.
These are different mathematical objects, and cryptographic problems
can use either one. In conventional elliptic-curve cryptography, one
usually fixes a large cyclic subgroup $\langle P\rangle$ of
$E(\mathbb F_q)$ and studies the recovery of $d$ from $P$ and
$[d]P$. In isogeny constructions, one studies maps between curves,
endomorphism rings, or actions of ideal class groups on suitable
families. The theory of orders speaks especially directly to the
second collection of objects. It also clarifies arithmetic used in
point counting and efficient scalar multiplication, but a description
of an endomorphism ring is not in itself a discrete-logarithm algorithm.

To keep these roles separate, write
$\operatorname{End}_{k}(E)$ for maps defined over $k$,
$\operatorname{End}_{\bar k}(E)$ for geometric endomorphisms,
and $\operatorname{End}^{0}_{k}(E)=
\operatorname{End}_{k}(E)\otimes_{\mathbb Z}\mathbb Q$.
Over finite fields, confusing the first two rings can turn a
quadratic order into a claimed quaternion order with the wrong
field of definition. Sutherland explicitly emphasizes this
distinction in his treatment of ordinary and supersingular curves
\cite{Sutherland2023Ordinary}.

\section{Complex tori and scalar endomorphisms}

Let $L=\mathbb Z\omega_1+\mathbb Z\omega_2\subseteq\mathbb C$
with $\omega_1,\omega_2$ linearly independent over $\mathbb R$.
The quotient $\mathbb C/L$ is a compact one-dimensional complex
Lie group. The analytic uniformization theorem identifies such
groups with complex elliptic curves; that theorem is an external
input here. The following description of their homomorphisms has
a short direct proof once the elementary lifting property of the
covering map $\mathbb C\to\mathbb C/L$ is available.

\begin{proposition}[Homomorphisms of complex tori]
\label{prop:torus-hom}
Every holomorphic group homomorphism
$f:\mathbb C/L\to\mathbb C/M$ is induced by a unique scalar
$\alpha\in\mathbb C$ satisfying $\alpha L\subseteq M$.
It is an isomorphism if and only if $\alpha L=M$.
\end{proposition}

\begin{proof}
Lift $f$ to a holomorphic map $F:\mathbb C\to\mathbb C$ with
$F(0)=0$. The group-homomorphism property implies
$F(z+w)-F(z)-F(w)\in M$. This is a continuous function of
$(z,w)$ taking values in the discrete set $M$, so it is constant;
at $(0,0)$ it is zero. Thus $F$ is additive. Differentiating
$F(z+w)=F(z)+F(w)$ with respect to $w$ at zero gives
$F'(z)=F'(0)=\alpha$. Therefore $F(z)=\alpha z$.
For $\lambda\in L$, the points $z$ and $z+\lambda$ have the
same image in the source torus, so $\alpha\lambda\in M$.
Conversely, that containment makes $z+L\mapsto\alpha z+M$
well-defined and holomorphic. A nonzero scalar has an inverse
on the quotient exactly when $\alpha L=M$. The lift fixing
zero is unique, so $\alpha$ is unique.
\end{proof}

In particular,
\[
 \operatorname{End}(\mathbb C/L)
   =O(L):=\{\alpha\in\mathbb C:\alpha L\subseteq L\}.
\]
This is the multiplier-order construction from the previous chapter,
now restricted to scalar multiplication on the plane. The
connection is literal: a holomorphic map of the torus is recorded
by an arithmetic operation preserving its period lattice.

\begin{theorem}[The elementary complex-multiplication dichotomy]
\label{thm:cm-dichotomy}
For a full lattice $L\subseteq\mathbb C$, either $O(L)=\mathbb Z$
or $O(L)$ is an order in an imaginary quadratic field.
\end{theorem}

\begin{proof}
Multiplication by $\alpha\in O(L)$ is represented in a
$\mathbb Z$-basis of $L$ by a matrix $A_\alpha\in M_2(\mathbb Z)$.
Its real-linear eigenvalues, after complexification, are
$\alpha$ and $\bar\alpha$. Consequently
$\alpha$ satisfies
$X^2-\operatorname{tr}(A_\alpha)X+\det A_\alpha=0$.
If $\alpha$ is real, multiplication on the underlying real plane
is $\alpha$ times the identity in every real basis. Its matrix
entries are integral, so $\alpha\in\mathbb Z$.

If all elements are real, we have obtained $O(L)=\mathbb Z$.
Otherwise choose nonreal $\alpha\in O(L)$. Its quadratic
polynomial is irreducible over $\mathbb Q$, since a real quadratic
polynomial with a rational root has both roots real. Put
$K=\mathbb Q(\alpha)$, an imaginary quadratic field. Every
$\beta\in O(L)$ commutes with $\alpha$, so $A_\beta$ lies in
the centralizer of $A_\alpha$ in $M_2(\mathbb Q)$. This
centralizer is $\mathbb Q[A_\alpha]$: choose any nonzero
$v\in\mathbb Q^2$; $v,A_\alpha v$ form a basis, and a
commuting operator is uniquely determined by its value on $v$.
Hence the centralizer has dimension two, and it already contains
the two-dimensional algebra $\mathbb Q[A_\alpha]$.

It follows that every $\beta$ belongs to $K$. The additive group
$O(L)$ embeds in $M_2(\mathbb Z)$ and is therefore finitely
generated. Its rational span is $K$, since it contains $1$ and
$\alpha$. Thus it is a full $\mathbb Z$-lattice and a ring
containing $1$, which is exactly an order in $K$.
\end{proof}

This argument proves an algebraic classification of scalar
endomorphisms. It does not prove that every finite-field elliptic
curve can be lifted with all of its geometric endomorphisms to
characteristic zero. In particular, a quaternion algebra cannot
embed into the scalar algebra $\mathbb C$. The supersingular
case below therefore needs a different geometric setting.

\section{Quadratic orders and the conductor}

Let $K$ be a quadratic number field with ring of integers
$\mathcal O_K=\mathbb Z+\mathbb Z\omega$. A subring finite over
$\mathbb Z$ consists of algebraic integers, so every order
$O\subseteq K$ is contained in $\mathcal O_K$. The following
description gives the integral datum that is lost on passing from
$O$ to $O\otimes\mathbb Q$.

\begin{theorem}[Quadratic conductor classification]
\label{thm:quadratic-conductor}
Every order in $K$ has the form
\[
 O_f=\mathbb Z+f\mathcal O_K
     =\mathbb Z+\mathbb Z(f\omega)
\]
for a unique integer $f\ge1$. Its conductor
$\{x\in\mathcal O_K:x\mathcal O_K\subseteq O_f\}$ equals
$f\mathcal O_K$, and
$\operatorname{disc}(O_f)=f^2\operatorname{disc}(\mathcal O_K)$.
\end{theorem}

\begin{proof}
Since $\mathbb Z\subseteq O$, the subgroup
$O/\mathbb Z$ of $\mathcal O_K/\mathbb Z\simeq\mathbb Z$
is $f\mathbb Z$ for a unique $f\ge1$; it is nonzero because
$O$ has rank two. Therefore $O=\mathbb Z+\mathbb Z(f\omega)$.
Conversely, if $\omega^2=t\omega-n$ with $t,n\in\mathbb Z$,
then
$(f\omega)^2=ft(f\omega)-f^2n$, proving closure under
multiplication for every $f$.

Certainly $f\mathcal O_K$ lies in the conductor. Write an element
of the conductor as $x=a+b\omega$. Since $x=x\cdot1\in O_f$,
we have $f\mid b$. Since
$x\omega=-bn+(a+bt)\omega$ also lies in $O_f$, we have
$f\mid a+bt$, hence $f\mid a$. Thus
$x\in f\mathcal O_K$, giving the reverse containment.
The change from the basis $(1,\omega)$ to $(1,f\omega)$ has
determinant $f$, so the discriminant statement follows from the
change-of-lattice formula proved in the previous chapter.
\end{proof}

\begin{example}[A noninvertible ideal in a nonmaximal order]
Let $O=\mathbb Z[2i]=\mathbb Z+2\mathbb Z i$ and
$I=2\mathbb Z[i]$. Then $I$ is an $O$-ideal and
$O(I)=\mathbb Z[i]$, which strictly contains $O$. If $I$ were
invertible as an $O$-ideal with inverse $J$, the containment
$xI\subseteq I$ would imply
$xO=xIJ\subseteq IJ=O$, hence $x\in O$. The element
$i\in O(I)\setminus O$ contradicts this. Thus the ideal is
not invertible over $O$, despite being a valid full
arithmetic lattice.
\end{example}

This example pinpoints a necessary qualification in class-group
arguments. One cannot form a group from all nonzero ideals of an
arbitrary order merely by declaring multiplication. The invertible
fractional ideals do form a group; the noninvertible ideals are
additional module-theoretic data.

\section{Ideal classes as a torsor of elliptic curves}

Let $O$ be an imaginary quadratic order. Denote by
$\operatorname{Pic}(O)$ the group of invertible fractional
$O$-ideals modulo principal fractional ideals. For an invertible
ideal $I$, its multiplier ring is $O$ by the argument in the
preceding example, so $\mathbb C/I$ has endomorphism ring $O$.
The usual complex-multiplication correspondence identifies the
proper ideal classes with the corresponding complex elliptic curves;
the general correspondence is an external theorem, while the
following action on invertible ideal lattices is elementary
\cite{Sutherland2023CMTorsor}.

\begin{proposition}[The ideal-class torsor]
On homothety classes of invertible fractional $O$-ideals, the rule
\[
 [\mathfrak a]\cdot[I]=[\mathfrak a^{-1}I]
\]
defines a free and transitive action of $\operatorname{Pic}(O)$.
\end{proposition}

\begin{proof}
Changing $\mathfrak a$ or $I$ by a nonzero scalar changes
$\mathfrak a^{-1}I$ by a nonzero scalar, so the rule is
well-defined. Multiplication of ideals is associative and
commutative, so it gives a group action. If
$[\mathfrak a^{-1}I]=[I]$, multiply the equality of ideal classes
by $[I]^{-1}$ to get $[\mathfrak a]=[O]$; thus the action is
free. Given $[I]$ and $[J]$, the class
$[\mathfrak a]=[I][J]^{-1}$ satisfies
$[\mathfrak a^{-1}I]=[J]$, proving transitivity.
\end{proof}

A torsor is a set with a free transitive group action, not a group
with a specified identity. Choosing a base curve identifies its
isomorphism classes with a class group abstractly. That choice
does not automatically make discrete logarithms, ideal recovery,
or group multiplication on curve encodings computationally easy.
In algorithmic discussions, this distinction separates an efficient
group action from an efficiently represented group.

For an integral invertible ideal $\mathfrak a\subseteq O$,
$I\subseteq\mathfrak a^{-1}I$, and the quotient map induces an
isogeny of complex tori
\[
 \phi_{\mathfrak a}:\mathbb C/I\longrightarrow
 \mathbb C/(\mathfrak a^{-1}I).
\]
Its kernel is exactly the finite group
$\mathfrak a^{-1}I/I$. This follows by testing which cosets
$z+I$ become zero in the target. For a principal ideal
$(\alpha)$, the kernel size is
$[I:\alpha I]$, equal to the absolute determinant of
multiplication by $\alpha$ and hence to
$|N_{K/\mathbb Q}(\alpha)|$. The standard general formula is
$\deg\phi_{\mathfrak a}=[O:\mathfrak a]$ for invertible
integral ideals; its extension from principal to invertible
ideals uses local freeness and multiplication of ideal norms
\cite{Sutherland2023ModularCurves}.

\begin{example}[A visible endomorphism of the square torus]
For $I=\mathbb Z[i]$, multiplication by $1+i$ has matrix
\[
 \begin{pmatrix}1&-1\\1&1\end{pmatrix}
\]
in the basis $(1,i)$. Its determinant is two. It therefore gives
an isogeny of degree two and a kernel of order two. The matrix
also shows that multiplication rotates by $\pi/4$ and scales
length by $\sqrt2$. This is an example where the integral
operation, the geometric norm, and the degree can all be compared
directly. The comparison works because the metric and embedding
have been explicitly supplied.
\end{example}

\section{Finite-field endomorphisms and Frobenius}

The geometric classification is a foundational external theorem:
in characteristic zero an elliptic curve has geometric endomorphism
ring $\mathbb Z$ or an imaginary quadratic order. Over an algebraic
closure of a finite field, an ordinary elliptic curve has an
imaginary quadratic endomorphism order, while a supersingular curve
has a maximal order in the rational quaternion algebra ramified at
$p$ and infinity. The degree is the reduced norm in the quaternion
case. These statements require the theory of isogenies and cannot
be deduced from the complex-torus argument alone
\cite[Theorem~42.1.9]{Voight2021Supersingular}.

If $E$ is defined over $\mathbb F_q$, its $q$-power Frobenius
$\pi$ satisfies
\[
 \pi^2-[t]\pi+[q]=0,
 \qquad t=q+1-\#E(\mathbb F_q),
 \qquad |t|\le2\sqrt q.
\]
The characteristic equation and Hasse's bound are external inputs
from elliptic-curve theory \cite{Sutherland2023Hasse,Sutherland2023EndAlgebras}.
Once they are available, several useful arithmetic conclusions
are immediate, and it is useful to isolate them from those inputs.

\begin{proposition}[The Frobenius order]
Suppose $E/\mathbb F_q$ is ordinary, and put
$K=\mathbb Q(\pi)$. Then
\[
 \mathbb Z[\pi]\subseteq\operatorname{End}_{\mathbb F_q}(E)
 \subseteq\operatorname{End}_{\overline{\mathbb F}_q}(E)
 \subseteq\mathcal O_K.
\]
Moreover, if an intervening order $O$ has index
$h=[O:\mathbb Z[\pi]]$, then
$t^2-4q=h^2\operatorname{disc}(O)$.
\end{proposition}

\begin{proof}
Frobenius and integer multiplications are defined over
$\mathbb F_q$, giving the first inclusion. Extending the field
of definition gives the second. The ordinary classification
identifies the geometric endomorphism algebra with an imaginary
quadratic field; it contains the nonrational Frobenius and is
therefore $K$. Endomorphisms are integral, giving the final
inclusion. The trace and norm of $\pi$ are $t$ and $q$.
The trace Gram matrix in the basis $(1,\pi)$ is
\[
 \begin{pmatrix}2&t\\t&t^2-2q\end{pmatrix},
\]
whose determinant is $t^2-4q$. The index-squared discriminant
formula now proves the assertion.
\end{proof}

This proposition provides an explicit bridge to the conductor
classification. It constrains which orders can contain the
Frobenius order. It does not determine which one is the full
endomorphism ring without further information. Nor does knowing
$\#E(\mathbb F_q)$ produce discrete logarithms in a large
prime-order subgroup. Point counting and discrete logarithms
ask different questions even though both use the same curve.

\section{Tate modules as arithmetic lattices}

Let $\ell\ne p$ be a prime. Over an algebraic closure,
$E[\ell^n]$ is abstractly isomorphic to
$(\mathbb Z/\ell^n\mathbb Z)^2$. The inverse limit under
multiplication by $\ell$ is the Tate module
\[
 T_\ell(E)=\varprojlim_n E[\ell^n],
\]
a free $\mathbb Z_\ell$-module of rank two. Endomorphisms act
on every torsion level and therefore on this inverse limit.
This is a direct meeting of the completed-lattice
theory in the first volume with elliptic-curve arithmetic.

\begin{lemma}[Faithfulness on prime-to-characteristic torsion]
If an endomorphism $\alpha$ annihilates $E[\ell^n]$ for every
$n\ge1$, then $\alpha=0$.
\end{lemma}

\begin{proof}
A nonzero endomorphism is a finite isogeny, so its kernel has at
most $\deg\alpha$ geometric points. But $E[\ell^n]$ has
$\ell^{2n}$ geometric points for every $n$. If all these groups
lie in the kernel, their cardinalities eventually exceed
$\deg\alpha$, a contradiction.
\end{proof}

Thus endomorphisms can be represented by compatible integral
matrices on Tate modules. Obtaining a full description of the
endomorphism ring from Galois-equivariant Tate-module maps is a
deeper theorem of Tate, not merely the faithfulness lemma above
\cite{Tate1966Endomorphisms}. The distinction mirrors the
difference between an injective map into a matrix algebra and
a theorem identifying its image.

For supersingular $E$ with all geometric endomorphisms defined,
the maximal order has $\mathbb Z$-rank four. Away from $p$, its
completed algebra is split, and its action on prime-power torsion
relates it to $M_2(\mathbb Z_\ell)$. At $p$, the algebra is a
division quaternion algebra. The exceptional prime is not a
minor implementation detail: the geometry of inseparable
isogenies and the local ramification require a different
description. This explains why replacing all local factors by
the same $2\times2$ matrix model loses arithmetic information.

\section{Quaternion ideals and the Deuring correspondence}

Fix a supersingular curve $E_0/\overline{\mathbb F}_p$ and put
$O=\operatorname{End}(E_0)$,
$B=O\otimes_{\mathbb Z}\mathbb Q$. The Deuring correspondence
relates supersingular elliptic curves to ideal lattices over $O$.
Its module convention is crucial. The group
$\operatorname{Hom}(E_0,E)$ is a \emph{right} $O$-module by
precomposition:
\[
 f\cdot\alpha=f\circ\alpha.
\]
Postcomposition by an endomorphism of $E$ acts on the left.
With suitable identifications in the common rational quaternion
space, the Hom group becomes a full right $O$-ideal. In this
convention the other endomorphism ring is its left multiplier
order. The categorical correspondence is an external theorem;
Voight's treatment develops its ideals and morphisms explicitly
\cite[\S42.3]{Voight2021Supersingular}.

The next elementary lemma makes the multiplier identification
transparent once the correspondence has supplied the ideal.

\begin{lemma}[Endomorphisms of a right ideal lattice]
\label{lem:right-ideal-end}
Let $O$ be an order in a quaternion algebra $B$, and let
$I\subseteq B$ be a full right $O$-ideal. Then
\[
 \operatorname{End}_O(I)\simeq O_L(I).
\]
The map sends $a\in O_L(I)$ to left multiplication by $a$.
\end{lemma}

\begin{proof}
Left multiplication preserves $I$ exactly for elements of
$O_L(I)$ and commutes with right multiplication by $O$, so
the map is well-defined. It is injective because $I$ spans $B$.
For surjectivity, extend $f\in\operatorname{End}_O(I)$ to a
$\mathbb Q$-linear map $f_{\mathbb Q}:B\to B$. Since $O$
spans $B$, the extension is right $B$-linear. A right
$B$-linear endomorphism of $B$ is left multiplication by
$a=f_{\mathbb Q}(1)$, because
$f_{\mathbb Q}(x)=f_{\mathbb Q}(1\cdot x)=ax$.
Its preservation of $I$ says $a\in O_L(I)$. Composition of
left multiplications agrees with multiplication in $O_L(I)$,
so this is a ring isomorphism.
\end{proof}

One must also distinguish ideal classes from conjugacy classes
of maximal orders. Several ideal classes may have isomorphic
multiplier orders. Recording only an order up to conjugacy can
forget which ideal class, and therefore which object in the
correspondence, was present. This is an arithmetic instance of
a general principle: a derived invariant can classify less
information than the object from which it was computed.

\section{The quaternion norm and the meaning of degree}

For an elliptic-curve endomorphism $\alpha$, the dual isogeny
$\widehat\alpha$ satisfies
$\widehat\alpha\alpha=\alpha\widehat\alpha=[\deg\alpha]$.
These identities make the degree behave like a norm. In the
quaternion endomorphism algebra, the dual isogeny is the standard
involution, and degree agrees with the reduced norm. For a
nonzero quasi-endomorphism $\alpha/n$, the norm is
$\deg(\alpha)/n^2$. Thus passage from integral maps to rational
quasi-isogenies changes the integrality, while retaining the
norm law.

\begin{proposition}[Degree multiplicativity from duality]
Assume the dual-isogeny identities and
$\widehat{\beta\alpha}=\widehat\alpha\widehat\beta$.
Then $\deg(\beta\alpha)=\deg\beta\deg\alpha$.
\end{proposition}

\begin{proof}
Compute
\[
 \widehat{\beta\alpha}(\beta\alpha)
 =\widehat\alpha\widehat\beta\beta\alpha
 =\widehat\alpha[\deg\beta]\alpha
 =[\deg\beta\deg\alpha].
\]
On the other hand, the dual-isogeny identity for $\beta\alpha$
makes the left side $[\deg(\beta\alpha)]$. The map
$\mathbb Z\to\operatorname{End}(E)$, $n\mapsto[n]$, is
injective, so the integers are equal.
\end{proof}

This proof illustrates the appropriate form of a conditional
bridge: its assumptions are already strong theorems of
elliptic-curve geometry, and the conclusion is a short algebraic
consequence. It would be misleading to present the short
calculation as a replacement for the construction of the dual
isogeny or the proof of its properties.

For numerical visualization, a positive definite quaternion
norm can display shells $\operatorname{nrd}(x)=c$ in four
real coordinates. A sequence of three-dimensional slices, or a
projection controlled by a fourth-coordinate slider, can reveal
how an ideal lattice meets those shells. Such a picture helps
explain why a norm equation is a Diophantine constraint. A
rendered shell is not evidence that an unshown integer point
exists, and a point drawn with floating-point error is not a
verified ideal element.

\section{What these bridges establish for cryptography}

The arithmetic theory now supplies several precise translations.
Complex multiplication translates scalar endomorphisms into
imaginary quadratic orders. Ideal multiplication translates into
actions and isogenies when the necessary invertibility hypotheses
hold. Supersingular endomorphisms translate into maximal
quaternion orders. Hom groups translate into integral bimodules.
Prime-to-characteristic torsion translates part of the theory
into compatible finite matrix modules.

Each translation has a different computational cost. An abstract
isomorphism, an efficiently computable isomorphism, and a reduction
preserving a promised distribution are three separate claims.
To apply the translations to cryptographic security, one must
specify which of them is actually available. In particular, an
algorithm given an explicit endomorphism basis need not solve
the problem of obtaining that basis from a curve equation. A
class-group action given an ideal factorization need not reveal
the factorization from the resulting curve. A statement about
all ideals is not a statement about the small ideals selected
by a protocol's distribution.

The most useful way to state a proposed new algorithm is therefore
as a tuple of mathematical data. The input type specifies a curve,
an order, a basis, a subgroup, or an ideal. The promise specifies
ordinary or supersingular, maximal or nonmaximal, and any supplied
orientation. The output specifies a scalar, a morphism, a basis,
or a witness to equivalence. The cost model counts bit operations
or fault-tolerant quantum resources. The proof then separates
correctness from cost and from any average-case assumption.

These requirements expose several invalid shortcuts. The fact
that an order is a lattice does not reduce arbitrary lattice
cryptography to an elliptic-curve problem. A four-dimensional
quaternion lattice is not the same family as a growing-dimensional
module lattice. A positive norm on an endomorphism algebra is not
a spectral gap for a quantum field theory. An exact functor
preserving module homology does not by itself supply a reversible
quantum circuit. Later chapters can use the constructions above,
but each claimed transfer must carry its own missing hypotheses.

\section{A worked comparison of three rank-two objects}

It is instructive to hold the rank fixed while changing the
structure. First consider $L_1=\mathbb Z^2$ as an abelian group.
Its automorphisms are $\mathrm{GL}_2(\mathbb Z)$, and it has no
preferred notion of length. Next equip the same module with the
standard dot product. Its isometries are the much smaller group
$\mathrm{GL}_2(\mathbb Z)\cap O(2)$, consisting of signed
permutation matrices. Finally identify it with
$\mathbb Z[i]\subseteq\mathbb C$ and remember multiplication
by $i$. Its module endomorphisms commuting with $i$ form the
ring $\mathbb Z[i]$.

\begin{proposition}[Centralizer of the quarter-turn]
Let $J=\left(\begin{smallmatrix}0&-1\\1&0\end{smallmatrix}\right)$.
An integral matrix commutes with $J$ if and only if it has the
form
$\left(\begin{smallmatrix}a&-b\\b&a\end{smallmatrix}\right)$
with $a,b\in\mathbb Z$. Its determinant is $a^2+b^2$.
\end{proposition}

\begin{proof}
Write the matrix as
$T=\left(\begin{smallmatrix}u&v\\w&z\end{smallmatrix}\right)$.
Equating entries in $TJ=JT$ gives $z=u$ and $v=-w$.
Set $a=u$, $b=w$. The determinant formula follows by direct
expansion. Conversely, any displayed matrix commutes with $J$
by multiplication.
\end{proof}

The determinant is the norm of $a+bi$, and the matrix is a
unit precisely when $a^2+b^2=1$. Hence the scalar automorphisms
of the square complex torus are $\{1,-1,i,-i\}$. The same
underlying rank-two group has therefore supported three different
automorphism questions. This small example is a useful check
whenever the monograph passes between algebraic, metric, and
complex-analytic structures.

\section{An exact arithmetic certificate for a quadratic order}

The order $O_f=\mathbb Z+f\mathbb Z i$ provides a complete
small certificate that can be checked without a numerical
approximation. In the basis $b_1=1$, $b_2=fi$, its multiplication
table is
\[
 b_1b_1=b_1,\quad b_1b_2=b_2b_1=b_2,
 \quad b_2b_2=-f^2b_1.
\]
Every structure constant is an integer, so the table certifies
closure under multiplication. The field trace sends
$a+bi$ to $2a$, and the trace matrix is
\[
 G_f=\begin{pmatrix}2&0\\0&-2f^2\end{pmatrix},
 \qquad \det G_f=-4f^2.
\]
Its negative determinant is not an error: the trace pairing
$\operatorname{Tr}(xy)$ is indefinite on this imaginary
quadratic field. The positive pairing
$\operatorname{Tr}(x\bar y)$ has matrix
$\operatorname{diag}(2,2f^2)$ instead. These two matrices encode
different bilinear structures on the same order.

\begin{proposition}[Trace dual of the Gaussian suborder]
For every integer $f\ge1$,
\[
 O_f^{\#}=\tfrac12\mathbb Z+\tfrac{i}{2f}\mathbb Z,
 \qquad [O_f^{\#}:O_f]=4f^2.
\]
\end{proposition}

\begin{proof}
Write $x=a+bi$ with $a,b\in\mathbb Q$. The conditions
$\operatorname{Tr}(x\cdot1)\in\mathbb Z$ and
$\operatorname{Tr}(x\cdot fi)\in\mathbb Z$ are
$2a\in\mathbb Z$ and $-2fb\in\mathbb Z$. They are
sufficient because $1,fi$ generate $O_f$ additively. This
proves the displayed description. In the first coordinate,
the quotient $\tfrac12\mathbb Z/\mathbb Z$ has size two.
In the imaginary coordinate,
$\tfrac1{2f}\mathbb Z/f\mathbb Z$ has size $2f^2$.
Multiplying their sizes gives $4f^2$.
\end{proof}

For $f=2$, the order has discriminant $-16$, its maximal
overorder $\mathbb Z[i]$ has discriminant $-4$, and the
index is two. The discriminant formula reads
$-16=2^2(-4)$. The conductor is $2\mathbb Z[i]$, which
also supplied the noninvertible ideal above. Its multiplier
ring is larger than $O_2$, so no inverse ideal over $O_2$
can exist. Thus multiplication, discriminant, conductor,
trace dual, and failure of invertibility are mutually
consistent in one exact example.

An implementation should output these integer matrices and
containment witnesses, not merely a decimal determinant.
A diagram may plot the corresponding point sets, but the
matrix identities are the certificates. The same separation
between exact witnesses and geometric display is useful in
larger quaternion calculations and in the interactive
companions to this monograph.

\section{A research program that is actually implied}

The preceding material supports a concrete, testable program.
One can represent orders and their ideals exactly by integral
bases, verify multiplication tables, compute trace matrices, and
check discriminant-index identities. One can formalize scalar
actions and centralizers in finite dimensions before attempting
the full geometric classification. One can compare localizations
of a specified lattice, keeping the difference between local
isomorphism and global equivalence visible. One can design
interactive diagrams whose point coordinates are generated from
verified arithmetic rather than artistic approximations.

The next mathematical advance would require a specific theorem
that the existing translations do not yet provide. Examples are
a distortion-controlled reduction between explicit arithmetic
norm problems, a sharper resource analysis for an exact quantum
arithmetic subroutine, or a new effective bound with fully stated
analytic assumptions. Each is a well-posed research target. None
follows simply by adjoining the names of order theory, elliptic
complexes, quantum algorithms, and Yang--Mills theory.

The didactic value of Roggenkamp and Huber-Dyson's framework is
that it forces the integral and categorical data to be stated
before any such claim is made. In that sense the books provide a
rigorous foundation for investigation: they clarify the objects
and identify obstructions that a successful new connection would
have to overcome.

